\section{Roots of a bond threshold}
\label{sec:h-timing}

A condition \(B_{ij}(t)=B_*\) is one real constraint on an almost-periodic
function whose frequencies are eigenvalue differences of \(H\).
For \(B(t)=B_0+C\cos(\omega t+\varphi)\) the crossings, when they exist, are
\[
t=\frac1\omega\Bigl(\pm\arccos\frac{B_*-B_0}{C}-\varphi\Bigr)+\frac{2\pi n}{\omega}.
\]
Equation~\eqref{eq:mudot} is two real constraints.
A generic one-dimensional orbit misses it, which is compatible with a
unitary flow that does not fire and with perturbations that stay small.
Several edges meeting a codimension-\(1\) level in the same window leave a
tie of the same kind as a degenerate re-pairing
(Section~\ref{sec:h-rewire}).
The level \(B_*\) still contains the two unfixed ratios of the conditional
conservation identity.
The crossing times of every candidate level are functions of the single
process \(B_{ij}(t)\).
Each crossing selects the radius in Proposition~\ref{prop:circle} and
leaves the angle.
The two unfixed ratios move that radius.
They do not pick the point.
Proposition~\ref{prop:Bdot} is the velocity of the same process:
\(\dot B_{ij}\) is a sum of pair currents on the other neighbours of the
two endpoints, and \(\mathcal{J}_{ij}\) itself does not appear.
On an isolated edge the sum is empty, \(\dot B_{ij}=0\) for every state,
and a level \(B_*\) is not crossed.
The current on that edge is the oscillator of Proposition~\ref{prop:Jdot},
at frequency \(2\), which is the spectral gap of \(K_2\).
The acceleration of the same process is Proposition~\ref{prop:Bddot}:
a sum of bonds along walks of length \(2\).
A real state has every current zero, so every bond is instantaneously
stationary, while that acceleration need not vanish.
On an isolated edge both the velocity and the acceleration of \(B\) vanish
for every state, and the quantity that moves with the angle is the product
in Proposition~\ref{prop:radius}.
